%%%%%%%%%%%%%%%%%%%%%%%%%%%%%%%%%%%%%%%%%%%%%%%%%%%%%%%%%%%%
%% Lecture 06
%%%%%%%%%%%%%%%%%%%%%%%%%%%%%%%%%%%%%%%%%%%%%%%%%%%%%%%%%%%%

\lecture
{05}
{Wednesday, 19 August 2026}
{Introduction to Probability}
{Dr. Nishant Chandgotia}

%%%%%%%%%%%%%%%%%%%%%%%%%%%%%%%%%%%%%%%%%%%%%%%%%%%%%%%%%%%%
\section{Continuous Random Variables}

%%%%%%%%%%%%%%%%%%%%%%%%%%%%%%%%%%%%%%%%%%%%%%%%%%%%%%%%%%%%

\begin{question}{Quiz-Question-Discussion}{quiz-question-discussion}
The last question from the quiz showed that pairwise independence and
mutual independence are different notions. What is the distinction
between the two?
\end{question}

%%%%%%%%%%%%%%%%%%%%%%%%%%%%%%%%%%%%%%%%%%%%%%%%%%%%%%%%%%%%

\subsection{Continuous Random Variables}

Let
\(\qquad 
X:(\Omega,\mathcal{F},\mathbb{P})
\longrightarrow
(S,\rho)
\)
be a random variable.

In many cases, the state space is one of
\(
\mathbb{R}, ~  ~ ~
\mathbb{R}^d,
\)
equipped with the Borel $\sigma$-algebra.

A random variable is called \emph{continuous} when its distribution
does not have atoms. In particular, for a real-valued continuous random
variable,
\(
\mathbb{P}(X=s)=0,
\qquad
\forall s\in\mathbb{R}.
\)

The distribution of $X$ is the push-forward of $\mathbb{P}$ via $X$:

%%%%%%%%%%%%%%%%%%%%%%%%%%%%%%%%%%%%%%%%%%%%%%%%%%%%%%%%%%%%
\subsection{Absolutely Continuous and Singular Distributions}

Let $\mu_X=X_*(\mathbb{P})$ denote the distribution of $X$.

\begin{definition}{Absolutely Continuous}{absolutely-continuous-definition}
The distribution $\mu_X$ is said to be \emph{absolutely continuous}
with respect to Lebesgue measure $\lambda$ if there exists a
non-negative function
\(
f\in L^1(\mathbb{R})
\)
such that
\(
\int_{\mathbb{R}}f\,d\lambda=1
\)
and
\(
\mu_X(A)
=
\mathbb{P}(X\in A)
=
\int_A f(x)\,d\lambda(x)
\)
for every Borel set $A\subseteq\mathbb{R}$.

The function $f$ is called the \emph{density} of $X$.
\end{definition}

\begin{remark}

The statement that the distribution of a random variable is absolutely
continuous means that its distribution is absolutely continuous with
respect to Lebesgue measure.

On the other hand, a distribution is called \emph{singular} with
respect to Lebesgue measure if it is concentrated on a set of Lebesgue
measure zero.

\end{remark}

More precisely, if
\(
\nu=X_*(\mathbb{P}),
\)
then $\nu$ is singular with respect to Lebesgue measure if there exists
a Borel set $A\subseteq\mathbb{R}$ such that

\[
\nu(A)=1,
\qquad
\lambda(A)=0.
\]

%%%%%%%%%%%%%%%%%%%%%%%%%%%%%%%%%%%%%%%%%%%%%%%%%%%%%%%%%%%%
\subsection{Examples of Continuous Random Variables}

\begin{enumerate}

    \item \textbf{Uniform random variable}

    The uniform distribution on $[0,1]$ has density

    \[
    f(x)
    =
    \begin{cases}
        1, & 0\leq x\leq1,\\
        0, & \text{otherwise}.
    \end{cases}
    \]

    \item \textbf{Normal random variable}

    Let

    \[
    X\sim N(\mu,\sigma^2),
    \qquad
    \sigma>0.
    \]

    Its density is

    \[
    f_X(x)
    =
    \frac{1}{\sqrt{2\pi}\sigma}
    \exp\left(
        -\frac{(x-\mu)^2}{2\sigma^2}
    \right).
    \]

    Hence, for $a<b$,

    \[
    \mathbb{P}(a\leq X\leq b)
    =
    \int_a^b
    \frac{1}{\sqrt{2\pi}\sigma}
    \exp\left(
        -\frac{(x-\mu)^2}{2\sigma^2}
    \right)
    dx.
    \]

    \item \textbf{Exponential random variable}

    Let

    \[
    \lambda>0,
    \qquad
    X\sim\operatorname{Exp}(\lambda).
    \]

    The random variable takes values in $[0,\infty)$ and has density

    \[
    f_X(x)
    =
    \lambda e^{-\lambda x},
    \qquad
    x\geq0.
    \]

    Thus, for $0\leq a<b$,

    \[
    \mathbb{P}(a<X<b)
    =
    \int_a^b
    \lambda e^{-\lambda x}\,dx.
    \]

    The exponential distribution has the \emph{memoryless property}:

    \[
    \mathbb{P}(X>s+t\mid X>s)
    =
    \mathbb{P}(X>t).
    \]

\end{enumerate}

%%%%%%%%%%%%%%%%%%%%%%%%%%%%%%%%%%%%%%%%%%%%%%%%%%%%%%%%%%%%
\subsection{A Singular Continuous Distribution}

Consider
\(
\Omega=\{0,1\}^{\mathbb{N}}
\)
with the product $\sigma$-algebra and the usual product probability
measure.

Define
\[
X(\omega)
=
\sum_{n=1}^{\infty}
\frac{\omega_n}{2^n}.
\]

Then

\[
X:\{0,1\}^{\mathbb{N}}\longrightarrow[0,1].
\]

\begin{question}{~}{ques-tion}
What is the distribution of $X$?
\end{question}

For a fixed sequence
\(
\widetilde{\omega}:
\{1,\ldots,n\}
\longrightarrow
\{0,1\},
\)
consider the corresponding dyadic interval

\[
I_{\widetilde{\omega}}
=
\left[
\sum_{i=1}^{n}
\frac{\widetilde{\omega}_i}{2^i},
\,
\sum_{i=1}^{n}
\frac{\widetilde{\omega}_i}{2^i}
+
\frac{1}{2^n}
\right].
\]

The event that the first $n$ coordinates agree with
$\widetilde{\omega}$ is the cylinder set
$[\widetilde{\omega}]$, and

\[
\mathbb{P}
\left(
X\in I_{\widetilde{\omega}}
\right)
=
\mathbb{P}
\left(
[\widetilde{\omega}]
\right)
=
2^{-n}.
\]

Since the length of the dyadic interval is also

\[
|I_{\widetilde{\omega}}|
=
2^{-n},
\]

we obtain, for dyadic intervals $I\subseteq[0,1]$,

\[
\mathbb{P}(X\in I)=|I|.
\]

This leads to the uniform distribution on $[0,1]$.

%%%%%%%%%%%%%%%%%%%%%%%%%%%%%%%%%%%%%%%%%%%%%%%%%%%%%%%%%%%%
\subsection{The Cantor Distribution}

Now consider

\[
\widetilde{X}(\omega)
=
\sum_{n=1}^{\infty}
\frac{2\omega_n}{3^n},
\qquad
\omega_n\in\{0,1\}.
\]

This is the ternary expansion obtained by using only the digits $0$
and $2$.

Consequently,

\[
\widetilde{X}(\omega)
\in C
\]

almost surely, where $C$ denotes the middle-thirds Cantor set.

Recall that

\[
\lambda(C)=0.
\]

Thus the distribution of $\widetilde{X}$ is concentrated on a set of
Lebesgue measure zero.

\begin{important}
The Cantor distribution is a singular continuous distribution.

It is continuous in the sense that

\[
\mathbb{P}(\widetilde{X}=\theta)=0
\]

for every $\theta\in\mathbb{R}$, but its distribution is singular with
respect to Lebesgue measure because it is concentrated on the Cantor
set.
\end{important}

To see why there are no atoms, consider the ternary expansion

\[
\theta
=
\sum_{n=1}^{\infty}
\frac{\theta_n}{3^n}.
\]

If some $\theta_i\notin\{0,2\}$, then

\[
\mathbb{P}(\widetilde{X}=\theta)=0.
\]

If $\theta_i\in\{0,2\}$ for every $i$, then fixing the first $n$
digits determines a cylinder set of probability
\(
2^{-n}.
\)

Therefore,

\(
\mathbb{P}(\widetilde{X}=\theta)
\leq
2^{-n}
\)
for every $n$, and hence
\(
\mathbb{P}(\widetilde{X}=\theta)=0.
\)

%%%%%%%%%%%%%%%%%%%%%%%%%%%%%%%%%%%%%%%%%%%%%%%%%%%%%%%%%%%%
\subsection{Bernoulli Convolution}

Consider
\(
\lambda\in\left(\frac12,1\right)
\)
and let
\(
\omega=(\omega_1,\omega_2,\ldots)
\in\{-1,1\}^{\mathbb{N}}.
\)

Define

\[
X_\lambda(\omega)
=
\sum_{n=1}^{\infty}
\omega_n\lambda^n.
\]

This is an example of a \emph{Bernoulli convolution}.

\begin{question}{~}{question-is-it-absolutely-cts}
For which values of $\lambda$ is the distribution of
$X_\lambda$ absolutely continuous?
\end{question}

This question is related to the classical problems surrounding
Bernoulli convolutions and the Erd\H{o}s conjecture.

\begin{question}{~}{What-is-the-distribution}
What distribution should we use to model a given random phenomenon?
\end{question}

%%%%%%%%%%%%%%%%%%%%%%%%%%%%%%%%%%%%%%%%%%%%%%%%%%%%%%%%%%%%
\section{Bertrand's Paradox}

\begin{question}{~}{what-is-the-probability}
Suppose we have a unit disc and an inscribed equilateral triangle.
Consider a uniformly chosen chord of the disc.

What is the probability that the chord has length greater than the
side length of the triangle?
\end{question}


Let $AB$ denote the side of the inscribed equilateral triangle and
$CD$ denote the randomly chosen chord. We want to determine

\[
\mathbb{P}(CD>AB).
\]

The difficulty is that the phrase \emph{uniformly chosen chord} does
not specify a unique probability model.

\begin{remark}

Bertrand's paradox illustrates that the answer depends on how
``uniformly'' is interpreted.

\end{remark}

%%%%%%%%%%%%%%%%%%%%%%%%%%%%%%%%%%%%%%%%%%%%%%%%%%%%%%%%%%%%
\subsection{Method 1: Choosing Two Endpoints}

Choose the two endpoints $C$ and $D$ independently and uniformly from
the circumference of the circle.

Under this model,

\[
\mathbb{P}(CD>AB)=\frac13.
\]

%%%%%%%%%%%%%%%%%%%%%%%%%%%%%%%%%%%%%%%%%%%%%%%%%%%%%%%%%%%%
\subsection{Method 2: Choosing the Midpoint}

Choose a radius uniformly and then choose a point $M$ uniformly along
that radius. Construct the unique chord whose midpoint is $M$.

Under this model,

\[
\mathbb{P}(CD>AB)=\frac12.
\]

The difference in the two answers demonstrates the ambiguity in the
phrase ``uniformly chosen chord.''

\begin{question}{~}{question-on-distribution}
What is the distribution of the chord length under each of these
probability models?
\end{question}

%%%%%%%%%%%%%%%%%%%%%%%%%%%%%%%%%%%%%%%%%%%%%%%%%%%%%%%%%%%%
\section{Distribution as a Push-Forward}

The distribution of a random variable is simply the push-forward of
the underlying probability measure.

Suppose

\[
X:
(\Omega,\mathcal{F},\mathbb{P})
\longrightarrow
(\mathbb{R},\mathcal{B}(\mathbb{R}))
\]

is a real-valued random variable.

Its distribution is

\[
\nu=X_*(\mathbb{P}).
\]

Therefore,

\[
\mathbb{P}(X\in[a,b])
=
\nu([a,b])
=
\mathbb{P}
\left(
X^{-1}([a,b])
\right).
\]

%%%%%%%%%%%%%%%%%%%%%%%%%%%%%%%%%%%%%%%%%%%%%%%%%%%%%%%%%%%%
\subsection{Change of Variables}

Suppose

\[
\Omega=[a,b]
\]

with Lebesgue probability measure, and suppose $X$ is a sufficiently
regular monotone function.

Then the distribution of $X$ can be obtained from the push-forward
measure.

If $X$ is invertible, then

\[
\mathbb{P}(X\in[c,d])
=
\mathbb{P}
\left(
X^{-1}([c,d])
\right).
\]

For a monotone differentiable function,

\[
X^{-1}(d)-X^{-1}(c)
=
\int_c^d
(X^{-1})'(t)\,dt.
\]

Using

\[
(X^{-1})'(t)
=
\frac{1}{X'(X^{-1}(t))},
\]

we obtain

\[
\mathbb{P}(X\in[c,d])
=
\int_c^d
\frac{1}
{X'(X^{-1}(t))}
\,dt,
\]

whenever the required regularity and monotonicity assumptions hold.

%%%%%%%%%%%%%%%%%%%%%%%%%%%%%%%%%%%%%%%%%%%%%%%%%%%%%%%%%%%%
\subsection{Example: Chord Length}

\textbf{Method 1.}

Let

\[
\Omega=[0,1]
\]

and define

\[
X(\theta)
=
\left|
e^{\pi i\theta}-1
\right|.
\]

Since

\[
\left|e^{i\alpha}-1\right|
=
2\left|\sin\frac{\alpha}{2}\right|,
\]

we have

\[
X(\theta)
=
2\sin\left(\frac{\pi\theta}{2}\right),
\qquad
0\leq\theta\leq1.
\]

Thus

\[
X:[0,1]\longrightarrow[0,2].
\]

Its inverse is

\[
X^{-1}(x)
=
\frac{2}{\pi}
\arcsin\left(\frac{x}{2}\right).
\]

Therefore,

\[
\frac{d}{dx}X^{-1}(x)
=
\frac{1}
{\pi\sqrt{1-\frac{x^2}{4}}}.
\]

Hence the density of $X$ is

\[
f_X(x)
=
\frac{1}
{\pi\sqrt{1-\frac{x^2}{4}}},
\qquad
0<x<2.
\]

%%%%%%%%%%%%%%%%%%%%%%%%%%%%%%%%%%%%%%%%%%%%%%%%%%%%%%%%%%%%

\textbf{Method 2.}

Alternatively, consider

\[
X:[0,1]\longrightarrow[0,2]
\]

defined by

\[
X(\omega)
=
2\sqrt{1-\omega^2}.
\]

Here,

\[
X^{-1}(x)
=
\sqrt{1-\frac{x^2}{4}},
\qquad
0\leq x\leq2.
\]

The corresponding density is obtained by differentiating the inverse
appropriately.


%%%%%%%%%%%%%%%%%%%%%%%%%%%%%%%%%%%%%%%%%%%%%%%%%%%%%%%%%%%%

\begin{important}

The central point of Bertrand's paradox is that specifying the word
``uniformly'' is not enough. A probability model must specify the
underlying sample space and probability measure.

Different choices of the probability space and probability measure can
produce different distributions for the same geometric quantity.

\end{important}

%%%%%%%%%%%%%%%%%%%%%%%%%%%%%%%%%%%%%%%%%%%%%%%%%%%%%%%%%%%%


\lectureend
